\documentclass[10pt,a4paper]{amsart}
\usepackage[margin=19mm]{geometry}
\usepackage{amsmath,amssymb,mathtools}
\usepackage[scr=rsfs,cal=euler]{mathalfa}
\usepackage{xfrac}
\usepackage[unicode,hidelinks,bookmarks=false]{hyperref}
\newcommand{\ie}{i.e.\ }
\newcommand{\cf}{cf.\ }
\newcommand{\resp}{respectively\ }
\newcommand{\Subsection}{\subsection}
\providecommand{\nobreakdash}{\nobreak\hbox{-}}
\setlength{\emergencystretch}{2em}
\multlinegap0pt
\usepackage{amssymb}
\usepackage{tikz}
\usepackage{enumerate}
\usepackage{float}
\newtheorem{lemma}{Lemma}
        \DeclareMathOperator{\sys}{sys}
\title{Higher cosystoles of matroids}
\author{James Dylan Douthitt, Elana Israel, Lee Kennard}
\date{Source: arXiv:2605.20409v1 (CC BY 4.0)}
\begin{document}
\maketitle

\noindent\emph{Source and licence.} Higher cosystoles of matroids. Version arXiv:2605.20409v1 (CC BY 4.0), distributed by arXiv under the Creative Commons Attribution 4.0 International licence, http://creativecommons.org/licenses/by/4.0/, as recorded in the arXiv licence metadata for this exact version and shown on the abstract page https://arxiv.org/abs/2605.20409v1. Full licence notice: \texttt{licenses/arxiv-2605.20409-CC-BY-4.0.txt}.\par\smallskip

\noindent\textit{Editorial scope.} Counted unit: the article's lemma graph (source0.tex lines 1090--1092). The article's complete original proof of it is delivered with it (source0.tex lines 1094--1107, one \texttt{proof} environment of 14 lines). No mathematics is written, repaired or re-derived: the statement and the proof are the article's own lines, in the article's own notation, and no retained line is substituted or re-flowed.  No further source-local context is needed: every cross-reference of every retained line resolves inside the two delivered spans.  Nothing was omitted from any retained span, so the package carries no omission record.  No retained line contains a citation, so the wrapper carries no bibliography and no bibliography entry.  No retained line contains a figure environment, an image-inclusion command or drawing code, so no image is delivered and the package declares no asset dependency.  Editorial formatting only, in addition: an AMS \texttt{amsart} preamble that restates the article's own declarations and macros used by the retained lines, namely the environments \texttt{lemma} on separate counters, together with the standard packages those lines need.  The retained environments print in this document as Lemma 1 in order of appearance, whereas the article prints the same items with the numbering of its own full document.  The complete native source of the article as distributed in the arXiv e-print for this exact version is bundled unchanged as \texttt{sources/200989/source0.tex}.
\begin{lemma}\label{lem:3sys graph}% [$3$-systole of graphic matroids]\label
For $d \geq 3$, we have $\sys_3^*(M(K_{d+1})) = \frac{6}{d+1}$.
\end{lemma}

\begin{proof}
Consider the $d + 1$ cocircuits given by vertex bonds. Note that each edge in $K_{d+1}$ is contained in exactly two of these cocircuits, and note that any three satisfy the non-inclusion property since $d \geq 3$. Therefore, if $\mu$ is any weight function, and if $C_1^*, \ldots, C_{d+1}^*$ denotes these cocircuits ordered by their $\mu$-weight, then we have
    \[\sys_3^*(M(K_{d+1}), \mu) \leq \sum_{i=1}^3 \mu(C_i^*) \leq \frac{3}{d+1} \sum_{i=1}^{d+1} \mu(C_i^*) = \frac{6}{d+1}.\]
Since this estimate holds for any $\mu$, the upper bound follows. 

To see that equality holds, consider the constant weight function $\mu_1$ given by $\mu_1(e) = \binom{d+1}{2}^{-1}$ for all edges $e$. By the edge connectivity of $K_{d+1}$, every cocircuit of $M(K_{d+1})$ contains at least $d$ elements. 
Hence any three cocircuits $C_i^*$ satisfy
%We may choose three cocircuits with exactly $d$ elements such that
    \[\mu(C_1^*) + \mu(C_2^*) + \mu(C_3^*) \geq 3d\binom{d+1}{2}^{-1} = \frac{6}{d+1}.\]
Therefore
    \[\sys_3^*(M(K_{d+1})) 
    \geq \sys_3^*(M(K_{d+1}), \mu_1) \geq  \frac{6}{d+1}.\] 
Combined with the upper bound, this completes the proof.
\end{proof}

\end{document}
